\section{Ámbito de aplicación de Muon}

\subsection{¿En qué parámetros se usa Muon?}

\begin{warningbox}{Muon no es un sustituto universal}
Muon solo se usa para optimizar \textbf{las matrices de parámetros 2D de las capas ocultas}. Los siguientes parámetros siguen usando el optimizador estándar (AdamW o SGD con momento):
\begin{itemize}
  \item La capa de Input Embedding
  \item La capa Output (LM Head)
  \item Los parámetros escalares (Scalar)
  \item Los parámetros vectoriales unidimensionales (como el $\gamma$ de RMSNorm)
\end{itemize}
\end{warningbox}

En el entrenamiento real, Muon se usa \textbf{en combinación} con AdamW/SGD: las matrices de las capas ocultas pasan por Muon y el resto pasa por el optimizador estándar.

\subsection{Consideración especial para la capa de salida}

Para la capa de salida (la capa de clasificación), queremos que las distintas características aprendidas contribuyan con pesos diferentes a la clasificación final. Si también se blanquea la capa de salida, se borra esta diferencia de escala significativa. Por eso Muon no es adecuado para la capa de salida.

\subsection{Resumen de la sección}

El ámbito de aplicación de Muon se limita estrictamente a los parámetros de las matrices 2D de las capas ocultas, y se usa de manera complementaria con el optimizador estándar.

